\Article{8}
When experiments of this kind are carried out with magnetic needles suspended from a vertical thread, the reaction exerted by torsion must be taken into account in the more precise experiments. Let us distinguish two horizontal diameters of the thread. One, $D$, is at the lower end, where the needle is attached, and is parallel to the magnetic axis of the needle. The other, $E$, is at the upper end, where the thread is fixed, and is parallel to $D$ when the thread is untwisted. Suppose that $E$ makes an angle $v$ with the magnetic meridian and that the magnetic axis, or $D$, makes an angle $u$ with it. Experience shows that the torsional force is at least approximately proportional to the angle $v-u$. Let us therefore consider the moment with which this force tends to make the angle $u$ equal to $v$. Its value is $=(v-u)\theta$. Now the moment of the terrestrial magnetic force tending to decrease the angle $u$ is $=TM\sin u$, so the condition for equilibrium is given by the equation $(v-u)\theta=TM\sin u$. The smaller $\theta$ is compared with $TM$, the more real solutions this equation admits. Here we are concerned only with small values of $u$, so we can safely use the equation $(v-u)\theta=TMu$ in its place, or $\frac vu=\frac{TM}{\theta}+1$. In our instruments the upper end of the thread is attached to a movable horizontal arm carrying a pointer that moves around the rim of a circle divided into degrees. Even if the collimation error, that is, the graduation corresponding to the value $v=0$, is still not known precisely enough, this pointer shows the difference between two values of $v$. Another part of the instrument likewise gives, with the greatest precision, the difference between the values of $u$ corresponding to equilibrium. It is clear that we obtain $\frac{TM}{\theta}+1$ by dividing the difference between two values of $v$ by the difference between the corresponding values of $u$ If a somewhat longer interval passes between the experiments made for this purpose, then, if the greatest precision is desired, we must take account of the daily variation of magnetic declination. This is easily done by making simultaneous observations with a second instrument, in which the upper end of the thread is left untouched. It hardly needs saying that the two instruments must be far enough apart that they cannot interfere perceptibly with each other.

\expandafter\def\csname HistoricalNoteText9\endcsname{Strictly speaking, these constituent threads are not truly single threads. They are merely of the kind merchants sell as unspun.}
\SourceLeaf{98}
\expandafter\def\csname HistoricalNoteText8\endcsname{The graduations on both scales increase from left to right.}
To show how precise observations of this kind can be, we give an example from the logbook. On $1832$ Sept $22$, leaving the collimation errors in place, we observed the following declinations $u$ and angles $v$\HistoricalNote{8}:
\begin{center}
\begin{tabular}{c|c||c|c||c}
&&\multicolumn{2}{c||}{First needle}&Second needle\\
Exp.&time&$u$&$v$&$u$\\\hline
I&$9^h33'$ a.m.&$+0^0\ 4'19''5$&$300^0$&$+0^0\ 2'12''1$\\
II&$9\ 57$&$-0\ 0\ 19{,}6$&$240$&$+0\ 1\ 37{,}7$\\
III&$10\ 16$&$-0\ 4\ 40{,}5$&$180$&$+0\ 1\ 18{,}8$
\end{tabular}
\end{center}
The declinations of the first needle, reduced to the state at the first observation, are therefore
\begin{center}
\begin{tabular}{crr}
I.&$u=+0^0\ 4'19''5$&$v=300^0$\\
II&$+0\ 0\ 14{,}8$&$240$\\
III&$-0\ 3\ 47{,}2$&$180$
\end{tabular}
\end{center}
Combining the observations gives the following values of the fraction $\frac{TM}{\theta}$
\begin{center}
\begin{tabular}{cr}
I and II $.\ .\ .\ .\ .$&$881{,}7$\\
II and III $.\ .\ .\ .\ .$&$891{,}5$\\
I and III $.\ .\ .\ .\ .$&$886{,}6$
\end{tabular}
\end{center}

The daily variations in magnetic declination are reduced by torsion in the ratio of unity to $\frac{n}{n+1}$, where we set $\frac{TM}{\theta}=n$. With threads whose torsion is as small as in the preceding example, this change can be regarded as imperceptible. As for the duration of the infinitely small oscillations, we can easily infer from the principles of dynamics that torsion reduces it in the ratio of unity to $\BareRoot{\frac{n}{n+1}}$. Strictly speaking, this applies to the case $v=0$. The formulas would hold generally if we set $\frac{TM\cos u^0}{\theta}=n$, with $u^0$ denoting the value of $u$ corresponding to equilibrium. But the difference will be quite imperceptible.

\Article{9}
The coefficient $\theta$ depends chiefly on the length, thickness and material of the thread. It also depends somewhat on temperature in metal threads and on the humidity of the air in silk threads. In the former, and perhaps in the latter too \textit{when they are single threads}, it appears to be \SourceLeaf{99}entirely independent of the weight they carry. Things are different with composite silk threads, such as those needed to support heavier needles. In these threads $\theta$ increases with the suspended weight, though it remains much smaller than the value of $\theta$ for a metal thread of the same length capable of bearing the same weight. For example, by a method entirely similar to the one we described in the preceding article, but using a different thread and needle, we found the value $n=597{,}4$ when the thread held the needle with only its usual fittings and the total weight was $496{,}2$ grams. We found $=424{,}8$ when the weight was increased to $710{,}8$ grams, so that in the first case $\theta=0{,}0016740\,TM$ and in the second $\theta=0{,}0023542\,TM$. The thread is $800$ millimeters long and is made up of $32$ single threads\HistoricalNote{9}. Each safely carries about $30$ grams, and they are arranged so that they are all under equal tension. It is probable, moreover, that the value of $\theta$ consists of a constant part and a part proportional to the weight, and that the constant part equals the sum of the values of $\theta$ for the individual single threads. Under this hypothesis, which experiments have yet to confirm sufficiently, the constant part in the example is found to be $=0{,}0001012\,TM$. The value of $\theta$ for a single thread is therefore $=0{,}00000316\,TM$. With the help of the value of $TM$ to be determined shortly, we infer from this hypothesis that the reaction of a single thread twisted through an arc equal to the radius $(57^0\ 18')$ is equivalent to the weight of a milligram pressing on a lever about $\frac1{17}$ millimeter long.

\Article{10}
If the oscillating body is a simple needle of regular shape and homogeneous mass, the moment of inertia $K$ can be calculated by known methods. If, for example, the body is a rectangular parallelepiped with sides $a,b,c$, density $=d$ and hence mass $q=abcd$, its moment of inertia about an axis through the center parallel to the side $c$ will be $=\frac1{12}(aa+bb)q$. In magnetic needles of this shape, the side $a$ parallel to the magnetic axis is usually much greater than the width $b$, so for rougher experiments it will even suffice to set $K=\frac1{12}aaq$. In more precise experiments, even when a simple needle is used, we would be reluctant to \SourceLeaf{100}admit the unsupported assumption of perfectly homogeneous mass and perfectly regular shape. In our experiments, where the oscillating needle is attached to more elaborate fittings, such a calculation is altogether impossible. We therefore had to devise another way to determine the moment $K$ with the greatest precision.

A transverse wooden bar was attached to the needle, and two equal weights hung from it, bearing on points $A,B$ of the bar through very sharp tips. These points lay on a horizontal line in the same vertical plane as the axis of suspension, at equal distances on either side of that axis. Denoting the mass of each weight by $p$ and the distance $AB$ by $2r$, the addition of this apparatus increases the moment $K$ by $C+2prr$, where $C$ is the sum of the moment of inertia of the bar about the line of suspension and the moments of inertia of the weights about the vertical axes through their tips and centers of gravity. Suppose, then, that we observe the oscillations both of the unloaded needle and of the needle loaded at two different distances, namely for $r=r'$ and $r=r''$. Suppose also that the periods of oscillation, reduced to infinitely small oscillations and corrected for the effect of torsion, are respectively found to be $=t,t',t''$. The three unknowns can then be obtained by combining the equations
\[
\begin{aligned}
TMtt&=\pi\pi K\\
TMt't'&=\pi\pi(K+C+2pr'r')\\
TMt''t''&=\pi\pi(K+C+2pr''r'')
\end{aligned}
\]
These unknowns are $TM,K$ and $C$. We will attain still greater precision if we observe the oscillations for several values of $r$, say $r=r',r'',r'''$ etc., with corresponding periods $t',t'',t'''$ etc., and determine the two unknowns $x,y$ by the method of least squares so that the following equations are satisfied as closely as possible
\[
\begin{aligned}
t'&=\BareRoot{\frac{r'r'+y}{x}}\\
t''&=\BareRoot{\frac{r''r''+y}{x}}\\
t'''&=\BareRoot{\frac{r'''r'''+y}{x}}\ \text{etc.}
\end{aligned}
\]
Having done this, we obtain
\[
\begin{aligned}
TM&=2\pi\pi px\\
K+C&=2py
\end{aligned}
\]

\SourceLeaf{101}
The following points about this method are also worth noting.

I. When the needle is not too highly polished, it suffices simply to set the wooden bar on it. If the surface of the needle is so smooth that friction cannot prevent the bar from sliding over it, the bar must be fastened more firmly to the rest of the apparatus so that the whole moves as a single rigid body. In either case, care must be taken that the points $A,B$ lie with sufficient precision on a horizontal straight line.

II. Since a series of such experiments takes several hours, the variation in the intensity of terrestrial magnetism over that interval cannot be neglected if the greatest precision is desired. Before proceeding to elimination, we must therefore reduce the observed periods to a constant value of $T$, for example to the mean value for the first experiment. This requires simultaneous observations on another needle, just as in Article 8. If they give the period of one oscillation at the mean times of the individual experiments as $=u,u',u'',u'''$ etc., we must use the following values in the calculation in place of the observed values $t',t'',t'''$ etc., respectively
\[
\frac{ut'}{u'},\qquad\frac{ut''}{u''},\qquad\frac{ut'''}{u'''}\quad\text{etc.}
\]

III. A similar caution applies to the variation of $M$ that arises from any change in temperature during the experiments. The reduction just described clearly includes this correction already if both needles undergo the same change in temperature and are affected by it in the same way.

IV. If we seek only the value of $TM$, the first experiment is clearly superfluous. Even so, it is useful to follow the experiments with the loaded needle at once with one using the unloaded needle. This also gives the value of $K$, which can be used for experiments made with the same needle at another time. Its value clearly remains constant, even though $T$ and $M$ may change over time.
